\subsection*{IV 20.1: Préliminaires et notations (cote 25, sites 4--8)}
The section introduces the projective bundle $P$ attached to a locally free Module of finite type, its dual projective bundle, the \emph{schéma des hyperplans}, and the hyperplane attached to an invertible quotient of the dual Module. The displayed pairing and its transpose lead to the defining ideal; a second description uses a local section, and a third identifies the hyperplane with a projective bundle. Grothendieck expressly prefers beginning with this third description (site 5). His references to paragraph 17 concern smoothness, and those to paragraph 19 concern \emph{diviseurs relatifs} and \emph{intersection complète transversale}; Chapter III is invoked for a proposed determination of the relative-divisor functor. The \emph{schéma d'incidence} and \emph{hyperplan universel} organize the inverse images $Y_{\xi}$ and the corresponding Modules in a diagram on site 7. The final paragraph on site 8 extends the term \emph{section hyperplane} from an immersion to the stated general morphism. Several visible slips or unsettled signs merit attention: site 4, formula (xx), prints $\underline{D}_P(1)$, followed by $\underline{L}_P(1)$ in the next sentence; the letter beneath the inverse-image discussion on site 7 has overprinted M/H traces, retained as a doubt. These forms have not been repaired from the mathematics. The handwriting also replaces \emph{linéaires} by \emph{coniques} in the opening title, adds \emph{fermé}, and revises the wording of the universal-hyperplane paragraph. Sources: \href{https://grothendieck.umontpellier.fr/web/viewer.html?file=../25.pdf\#page=5}{sites 4--8}, especially sites 5--7 for the reorganization and universal construction.
